% SGA 3，报告 VII_B，英文重建本——第 1 轮（中文译本）。
% 范围：依据原文件第 17--20 页，从第 1 节至 1.2.2 结束；止于原文件第 20 页的 1.2.3 之前。
% 合并阅读版页码： 530--533. 再版印刷页码： 519--522.
% 本范围内的原版页边定位： 492--495.
% 本范围内的编者注： 37--44.
% 依据： Polo--Gille Exp7B-13oct24.pdf,
% SHA-256 D303690E0195654BBE7C3DAA5D89EA23E24CD709A3575217234AC19AE8836694.
% 曾参阅 Jacob Reinhold 的 CC BY 4.0 Markdown，作为对比版本来源。
% 1.2.C 中的原生双箭头图是与源文绑定的 PDF 裁片；
% 其他所有陈列数学均为原生 LaTeX。

% 依据： Exp7B 原文件页码 17; 合并阅读版页码 530; 印刷页码 519.
\section*{1. 伪紧环上的形式簇}
\phantomsection
\label{sga3:VIIB:1}
\addcontentsline{toc}{section}{1. 伪紧环上的形式簇}

\medskip
\noindent\textbf{1.1.}
\phantomsection
\label{sga3:VIIB:1.1}
可以如下将一个形式概形（EGA~I，10.4.2）对应于每个伪紧环 \(A\)。
其底拓扑空间是 \(A\) 的开素理想（因而是极大理想）所成的集合 \(\Upsilon(A)\)，
赋予离散拓扑。对于每个子集 \(E\subset\Upsilon(A)\)，结构层在 \(E\) 上的截面空间是笛卡尔积
\[
  \prod_{\mathfrak m\in E}A_{\mathfrak m}.
\]
% 依据： Exp7B 原文件页码 18; 合并阅读版页码 531; 印刷页码 520.
如此得到的形式概形记为 \(\operatorname{Spf}(A)\)，称为 \(A\) 的
\emph{形式谱}。\footnote{\label{sga3:VIIB:nde:37}%
参见\ \SGARefLink{sga3:VIIB:0.1.2}{注记 0.1.2}。}

若 \(A\) 和 \(B\) 是两个伪紧环，则从 \(\operatorname{Spf}(B)\) 到
\(\operatorname{Spf}(A)\) 的\emph{态射}由一个映射
\[
  f:\Upsilon(B)\longrightarrow\Upsilon(A)
\]
和一族连续同态
\[
  f_y^\natural:A_{f(y)}\longrightarrow B_y,
  \qquad y\in\Upsilon(B).
\]
组成。这样的态射诱导连续同态
\[
  f^\natural:
  A=\prod_{x\in\Upsilon(A)}A_x
  \longrightarrow
  B=\prod_{y\in\Upsilon(B)}B_y.
\]
逆命题也成立：

\medskip
\noindent\textbf{命题。}
\phantomsection
\label{sga3:VIIB:1.1-spf-fully-faithful}
\textit{反变函子 \(A\mapsto\operatorname{Spf}(A)\) 是全忠实的。}

事实上，设 \(\phi:A\to B\) 是连续同态。\(B\) 的开极大理想的逆像
\(\phi^{-1}(\mathfrak n)\) 是 \(A\) 的开素理想，从而是 \(A\) 的极大理想。
因此，得到从 \(\Upsilon(B)\) 到 \(\Upsilon(A)\) 的映射
\[
  \mathfrak n\longmapsto\phi^{-1}(\mathfrak n)
\]
且 \(\phi\) 诱导连续同态
\[
  A_{\phi^{-1}(\mathfrak n)}\longrightarrow B_{\mathfrak n}.
\]
所以 \(\phi\) 诱导形式概形态射
\[
  \operatorname{Spf}(\phi):
  \operatorname{Spf}(B)\longrightarrow\operatorname{Spf}(A).
\]
容易验证
\[
  \operatorname{Spf}(\phi)^\natural=\phi,
  \qquad
  \operatorname{Spf}(f^\natural)=f
\]
对每个态射 \(f:\operatorname{Spf}(B)\to\operatorname{Spf}(A)\) 成立，这就证明了命题。

尽管这里仅讨论形如 \(\operatorname{Spf}(A)\) 的形式概形，我们仍将使用形式概形的语言，
以便将我们的论断建立在几何直观上。

\medskip
\noindent\textbf{1.2.}
\phantomsection
\label{sga3:VIIB:1.2}
设 \(k\) 是伪紧环。

\medskip
\noindent\textbf{定义 1.2.A。}\footnote{%
\label{sga3:VIIB:nde:38}编者添加了编号 \SGARefLink{sga3:VIIB:1.2.A}{1.2.A}--\SGARefLink{sga3:VIIB:1.2.E}{1.2.E}，以供下文引用。}
\phantomsection
\label{sga3:VIIB:1.2.A}
一个\emph{\(k\) 上的形式簇}是 \(\operatorname{Spf}(k)\) 上的形式概形 \(X\)，
它同构于某个预有限 \(k\)-代数 \(A\) 的形式 \(k\)-概形 \(\operatorname{Spf}(A)\)。
这时，代数 \(A\) 同构于 \(X\) 的\emph{仿射代数}，即 \(X\) 的结构层
\(\mathcal O_X\) 的截面代数。

以 \(\operatorname{Vaf}/k\) 表示 \(\operatorname{Spf}(k)\) 上的形式概形范畴中，
由形式 \(k\)-簇作为对象的满子范畴。\footnote{%
\label{sga3:VIIB:nde:39}由 \SGARefLink{sga3:VIIB:1.1}{1.1} 中态射的定义，
每个 \(X\in\operatorname{Ob}(\operatorname{Vaf}/k)\) 都是各单点形式簇
\(\operatorname{Spf}(\mathcal O_{X,x})\)（\(x\in X\)）在
\(\operatorname{Vaf}/k\) 中的直和。}

由 \SGARefLink{sga3:VIIB:1.1}{1.1}，函子 \(A\mapsto\operatorname{Spf}(A)\)
是从 \(\operatorname{Alp}/k\)（\SGARefLink{sga3:VIIB:0.4.1}{0.4.1}）到
\(\operatorname{Vaf}/k\) 的反等价。因此，由 \SGARefLink{sga3:VIIB:0.4.2}{0.4.2} 的推论：

\medskip
\noindent\textbf{命题 1.2.B。}
\phantomsection
\label{sga3:VIIB:1.2.B}
\textit{范畴 \(\operatorname{Vaf}/k\) 有投射极限和直极限。}\footnote{
\label{sga3:VIIB:nde:40}特别地，\(\operatorname{Vaf}/k\) 中存在余核；见下文。
此外，\(\operatorname{Vaf}/k\) 中过渡态射全为满射的滤过投射极限也可能为空；
参见\ 编者注~(34)。}

例如，设 \(X\in\operatorname{Ob}(\operatorname{Vaf}/k)\)，设
\[
  f:Y\longrightarrow X,
  \qquad
  g:Z\longrightarrow X
\]
是 \(X\) 上的两个形式 \(k\)-簇，且 \(A,B,C\) 分别是 \(X,Y,Z\) 的仿射代数。
% 依据：Exp7B 原文件第 19 页；合并阅读版第 532 页；印刷页码 521。
由 \SGARefLink{sga3:VIIB:0.4.1}{0.4.1}，纤维积 \(Y\times_X Z\) 的仿射代数等同于
\(B\widehat{\otimes}_A C\)，\footnote{%
\label{sga3:VIIB:nde:41}形式簇 \(Y\times_X Z\) 是各形式簇
\(B_x\widehat{\otimes}_{\mathcal O_{X,x}}C_x\) 对 \(x\in X\) 的直和，
其中 \(B_x\) 是满足 \(f(y)=x\) 的 \(y\in Y\) 所对应的
\(\mathcal O_{Y,y}\) 的积，\(C_x\) 类似定义。这将在
\SGARefLink{sga3:VIIB:2.5.1:corollary}{2.5.1} 中用到。}
因此，\(\operatorname{Vaf}/k\) 到所有形式 \(k\)-概形范畴的包含函子与有限投射极限交换
（参见\ EGA~I，10.7）。

形式 \(k\)-簇的直极限对应于其仿射代数的投射极限。

\medskip
\noindent\textbf{例 1.2.C（余核）。}
\phantomsection
\label{sga3:VIIB:1.2.C}
设 \(d,e:X\rightrightarrows Y\) 是 \(\operatorname{Vaf}/k\) 中的双箭头。
\(\operatorname{Coker}(d,e)\) 的仿射代数同构于在 \(X\) 和 \(Y\) 的仿射代数上
诱导的同态的核。\(\operatorname{Coker}(d,e)\) 还可以如下构造。
其底拓扑空间是底空间的余核。\footnote{%
\label{sga3:VIIB:nde:42}即 \(Y\) 关于对每个 \(x\in X\) 作认同
\(d(x)=e(x)\) 所得的商集，赋予商拓扑，此处该拓扑是离散的。}
设 \(p\) 是从 \(Y\) 的底集到此余核的典范投影，\(z\) 是余核的一点。
\(\operatorname{Coker}(d,e)\) 在 \(z\) 处的局部代数是以下双箭头的核
\begin{center}
  \phantomsection
  \label{sga3:VIIB:1.2.C-cokernel-double-arrow}
  % SGA3-VII-DIAG-084
  % SGA3-VII-DIAG-084
% Authority: Exp7B-13oct24.pdf, local p. 19.
% Crop: figures/Exp7B_localp019_example12C_cokernel_double_arrow.png
% Locator: tex/components/23_expose_VIIB_sections11_122_bicoalgebras_en.tex, Loop 1 line 183.
\(
d^\natural,e^\natural:\quad
\prod_{p(y)=z}\mathcal O_{Y,y}
\mathrel{\substack{\longrightarrow\\[-.65ex]\longrightarrow}}
\prod_{q(x)=z}\mathcal O_{X,x}
\)

\end{center}
其中
\[
  q=p\circ d=p\circ e,
\]
且 \(d^\natural\) 和 \(e^\natural\) 由同态
\(d_x^\natural\) 和 \(e_x^\natural\) 诱导，所用记号见 \SGARefLink{sga3:VIIB:1.1}{1.1}。

\medskip
\noindent\textbf{定义 1.2.D。}
\phantomsection
\label{sga3:VIIB:1.2.D}
设 \(\phi:k\to\ell\) 是伪紧环同态，\(X\) 是仿射代数为 \(A\) 的形式
\(k\)-簇。通过基变换得到的形式概形
\[
  X\times_{\operatorname{Spf}(k)}\operatorname{Spf}(\ell),
\]
是形式 \(\ell\)-簇。我们也将它记为
\[
  X\widehat{\otimes}_k\ell.
\]
其仿射代数是完备化张量积 \(A\widehat{\otimes}_k\ell\)
（参见\ 0.5 和 EGA~I，\S~10）。

\medskip
\noindent\textbf{注记 1.2.E。}
\phantomsection
\label{sga3:VIIB:1.2.E}
由于 \(k\) 上的每个形式簇都分解为 \(k\) 的局部分量上的形式簇，
\emph{在某些证明中我们将假定 \(k\) 是局部伪紧环}。

现在给出一些例子，同时确定我们的术语。

\medskip
\noindent\textbf{1.2.1.}
\phantomsection
\label{sga3:VIIB:1.2.1}
一个 \(k\)-\emph{函子}，按定义是从 \(\operatorname{Alf}/k\) 到
\((\mathbf{Sets})\) 的协变函子。由 \SGARefLink{sga3:VIIB:1.1}{1.1} 和
\SGARefLink{sga3:VIIB:0.4.2}{0.4.2}，我们可以通过将每个形式 \(k\)-簇 \(X\)
对应于如下函子，把
\[
  \operatorname{Vaf}/k
  \simeq
  (\operatorname{Alp}/k)^\circ
\]
等同于 \(k\)-函子范畴的一个满子范畴：
\[
  \operatorname{Alf}/k\longrightarrow(\mathbf{Sets}),
  \qquad
  C\longmapsto
  X(C)=
  \operatorname{Hom}_{\operatorname{Vaf}/k}
    \bigl(\operatorname{Spf}(C),X\bigr).
\]

% 依据： Exp7B 原文件页码 20; 合并阅读版页码 533; 印刷页码 522.
下文将遇到这样的 \(k\)-函子 \(\mathbf F\)：它将
\(\operatorname{Alf}/k\) 的每个对象 \(C\) 对应于一个 \(C\)-模 \(\mathbf F(C)\)，
并将 \(\operatorname{Alf}/k\) 的每个态射 \(\phi:C\to D\) 对应于一个 \(k\)-线性映射
\[
  \mathbf F(\phi):\mathbf F(C)\longrightarrow\mathbf F(D)
\]
使得对 \(x\in\mathbf F(C)\) 及 \(\lambda\in C\)，
\[
  \mathbf F(\phi)(\lambda x)
  =
  \phi(\lambda)\mathbf F(\phi)(x).
\]
由\SGARefLink{sga3:i:definition:3-1}{报告 I，3.1}，这样的 \(\mathbf F\)
配备一个 \(\mathbf O_k\)-模结构，其中 \(\mathbf O_k\) 表示将
\(\operatorname{Alf}/k\) 的每个对象 \(C\) 对应于其底环的环值 \(k\)-函子。

\medskip
\noindent\textbf{定义。}
\phantomsection
\label{sga3:VIIB:1.2.1-definitions}

\noindent\textit{\SGARefLink{sga3:vii:paragraph:5-1:equation:i}{(i)}}
若 \(\operatorname{Alf}/k\) 的每个态射 \(\phi:C\to D\) 都诱导从
\(D\otimes_C\mathbf F(C)\) 到 \(\mathbf F(D)\) 的双射，则称
\(\mathbf O_k\)-模 \(\mathbf F\) 为\emph{容许}的。\footnote{%
\label{sga3:VIIB:nde:43}编者引入了这个术语；参见\ \SGARefLink{sga3:VIIB:1.2.3}{1.2.3}
的编者注~(47)。}

\smallskip
\noindent\textit{\SGARefLink{sga3:vii:paragraph:5-1:equation:ii}{(ii)}}
若模 \(\mathbf F\) 是容许的，且对 \(\operatorname{Alf}/k\) 的每个对象
\(C\)，\(C\)-模 \(\mathbf F(C)\) 都平坦，则称其为\emph{平坦}的。\footnote{%
\label{sga3:VIIB:nde:44}由于 \(C\) 是 Artin 环，它因而是投射模。}

例如，若 \(M\) 是一个不一定伪紧的 \(k\)-模，我们如同
\SGARefLink{sga3:i:definition:4-6-1}{报告 I，4.6.1} 那样，用 \(\mathbf W(M)\)
表示 \(\mathbf O_k\)-模
\[
  C\longmapsto C\otimes_kM.
\]
于是，只要 \(M\) 在 \(k\) 上平坦，\(\mathbf W(M)\) 就平坦。

此外，函子 \(M\mapsto\mathbf W(M)\) 的右伴随是将每个
\(\mathbf O_k\)-模 \(\mathbf F\) 对应于如下 \(k\)-模的函子
\[
  \varprojlim_{\mathfrak l}\mathbf F(k/\mathfrak l),
\]
其中 \(\mathfrak l\) 遍历 \(k\) 的开理想。

\medskip
\noindent\textbf{1.2.2.}
\phantomsection
\label{sga3:VIIB:1.2.2}
下文总用诸如 \(\mathbf F\) 的粗体字母表示 \(\mathbf O_k\)-模。
当 \(k\) 是 Artin 环时，我们将 \(\mathbf F(k)\) 简记为 \(F\)。
这时，显然函子
\[
  \mathbf F\longmapsto F
\]
诱导从平坦 \(\mathbf O_k\)-模范畴到平坦 \(k\)-模范畴的等价！
采用这个术语，仅仅是为了使我们能够“仿佛 \(k\) 总是 Artin 环一样”进行推理。

按照报告 I，\S~\SGARefLink{sga3:i:definition:3-1}{3.1}，对于其他代数结构也采用类似约定。
因此，一个 \(\mathbf O_k\)-代数，或分别为 \(\mathbf O_k\)-余代数、
\(\mathbf O_k\)-Lie 代数、\(\mathbf O_k\)-\(p\)-Lie 代数，是由一个
\(\mathbf O_k\)-模 \(\mathbf M\) 及 \(\operatorname{Alf}/k\) 的每个对象 \(C\) 所对应的
\(\mathbf M(C)\) 上的 \(C\)-代数结构，或分别为 \(C\)-余代数、\(C\)-Lie 代数、
\(C\)-\(p\)-Lie 代数结构组成。我们还要求，对 \(\operatorname{Alf}/k\)
中的每个态射 \(\phi:C\to D\)，由 \(\mathbf M(\phi)\) 诱导的映射
\[
  D\otimes_C\mathbf M(C)\longrightarrow\mathbf M(D)
\]
是 \(D\)-代数同态，或分别为 \(D\)-余代数同态，等等。

最后，若 \(\mathbf F\) 和 \(\mathbf G\) 是两个 \(\mathbf O_k\)-模，
则 \(\mathbf F\otimes_k\mathbf G\) 表示 \(\mathbf O_k\)-模
\[
  C\longmapsto\mathbf F(C)\otimes_C\mathbf G(C).
\]

% 下一断点： Exp7B 原文件页码 20, 合并阅读版页码 533,
% 再版印刷页码 522, 1.2.3 及编者注 (45) 处，
% 以“先从下列引理开始”起。
